% Technical manuscript with native-image crops and JSON-derived predictive figures.
% Reconstruction base (additional evidence identified by the revision manifest):
% 90fbf8f89ac2634a9dbd2b27f81fc80d68b9943c (2026-11-11).
\documentclass[11pt,letterpaper]{article}
\usepackage[margin=0.92in,headheight=23pt,footskip=30pt]{geometry}
% Native OpenType fonts for XeLaTeX/LuaLaTeX; traditional newtx for pdfLaTeX.
% Avoid newtx virtual-font dependencies in xdvipdfmx on Windows TeX Live.
\usepackage{iftex}
\ifPDFTeX
  \usepackage[T1]{fontenc}
  \usepackage[utf8]{inputenc}
  \usepackage{newtxtext,newtxmath}
  \usepackage[final]{microtype}
  \usepackage{amsmath,mathtools}
\else
  \usepackage{amsmath,mathtools}
  \usepackage{fontspec}
  \usepackage{unicode-math}
  \setmainfont{TeXGyreTermesX-Regular.otf}[BoldFont=TeXGyreTermesX-Bold.otf, ItalicFont=TeXGyreTermesX-Italic.otf, BoldItalicFont=TeXGyreTermesX-BoldItalic.otf]
  \setmathfont{latinmodern-math.otf}
\fi
\usepackage{booktabs,tabularx,threeparttable,array,multirow}
\usepackage{siunitx}
\usepackage{xcolor}
\usepackage{tikz}
\usetikzlibrary{arrows.meta,calc,patterns,decorations.pathreplacing}
\usepackage{pgfplots}
\pgfplotsset{compat=1.18}
\usepackage{graphicx}
\usepackage{caption,subcaption}
\usepackage{fancyhdr}
\usepackage{titlesec}
\usepackage{enumitem}
\usepackage{hyperref}
\usepackage{cleveref}
\usepackage{listings}
\usepackage{url}
\usepackage{float}
\usepackage{longtable}
\usepackage{needspace}
\usepackage{etoolbox}
\usepackage{setspace}
\pretocmd{\section}{\Needspace{6\baselineskip}}{}{}
\pretocmd{\subsection}{\Needspace{6\baselineskip}}{}{}
\definecolor{ink}{HTML}{202934}
\definecolor{slate}{HTML}{495B68}
\definecolor{accent}{HTML}{287D84}
\definecolor{pale}{HTML}{E8F0F0}
\definecolor{fine}{HTML}{C9D3D7}
\definecolor{warm}{HTML}{6A5A45}
\hypersetup{colorlinks=true,linkcolor=accent,citecolor=accent,urlcolor=accent,pdftitle={Subpixel Geometric Reconstruction of a Cylindrical Candidate in Curiosity Mastcam-100 Imagery},pdfauthor={Justin Bogner}}
\titleformat{\section}{\color{ink}\large\bfseries}{\thesection.}{0.65em}{}
\titleformat{\subsection}{\color{ink}\normalsize\bfseries}{\thesubsection}{0.65em}{}
\titleformat{\subsubsection}{\color{slate}\normalsize\itshape}{\thesubsubsection}{0.65em}{}
\titlespacing*{\section}{0pt}{20pt}{8pt}
\titlespacing*{\subsection}{0pt}{13pt}{5pt}
\setlength{\parindent}{1.2em}
\setlength{\parskip}{2.5pt}
\setlength{\columnsep}{19pt}
\emergencystretch=2em
\linespread{1.075}
\captionsetup{font=small,labelfont={bf,color=ink},textfont={color=slate},skip=5pt}
\renewcommand\arraystretch{1.15}
\pagestyle{fancy}
\fancyhf{}
\fancyhead[L]{\footnotesize\textsc{Sol 3556 | Mastcam-100}}
\fancyhead[R]{\footnotesize\textsc{Geometric reconstruction}}
\fancyfoot[L]{\footnotesize Bogner \quad | \quad 13 November 2026}
\fancyfoot[R]{\footnotesize\thepage}
\renewcommand{\headrulewidth}{0.4pt}
\renewcommand{\headrule}{\hbox to\headwidth{\color{fine}\leaders\hrule height \headrulewidth\hfill}}
\newcommand{\rms}{\operatorname{RMS}}
\newcommand{\aicc}{\mathrm{AIC}_{\!c}}
\newcommand{\dif}{\mathop{}\!\mathrm{d}}
\newcommand{\px}{\,\mathrm{px}}
\newcommand{\DOF}{\mathrm{DOF}}
\newcommand{\sourcepath}[1]{\nolinkurl{#1}}
\newcommand{\mat}[1]{\mathbf{#1}}
\newcommand{\repo}{\href{https://github.com/CambrianMinds/martian-cylinder}{CambrianMinds/martian-cylinder}}
\newcommand{\pdsid}{\texttt{3556MR1025170721700585C00}}
\lstdefinestyle{code}{basicstyle=\ttfamily\footnotesize,frame=single,rulecolor=\color{fine},backgroundcolor=\color{pale!45},columns=fullflexible,breaklines=true,showstringspaces=false,aboveskip=9pt,belowskip=9pt}
\lstset{style=code}

\input{paper-analysis-values.tex}
\begin{document}
\thispagestyle{empty}
\vspace*{-0.24in}
\noindent{\color{accent}\rule{\linewidth}{2.2pt}}\par\vspace{13pt}
{\footnotesize\sffamily\bfseries\color{accent} QUANTITATIVE PLANETARY IMAGE ANALYSIS \quad / \quad RESEARCH ARTICLE}\\[12pt]
{\fontsize{22}{26}\selectfont\bfseries\color{ink} Subpixel Geometric Reconstruction of a Cylindrical Candidate in Curiosity Mastcam-100 Imagery}\\[8pt]
{\large\color{slate} Boundary metrology, frozen-shape prediction, and geological comparison}\\[17pt]
{\large Justin Bogner}\\[3pt]
{\small\color{slate} Cambrian Minds Media \quad | \quad 13 November 2026}\\[9pt]
{\footnotesize\color{slate} Data and scripts: \repo}\\[2pt]
{\footnotesize\color{slate} Curiosity Sol 3556 \quad / \quad \pdsid}\\[8pt]
\par\noindent{\color{fine}\rule{\linewidth}{0.5pt}}\par

\begin{abstract}
\noindent A centimeter-scale object in Curiosity Mastcam-100 imagery acquired on sol 3556 exhibits two curved terminations and a straight upper exterior boundary. Physically directed subpixel localization yields 35 near-rim, 24 far-termination, and 26 upper-boundary stations. A six-parameter common-section, common-axis model with exact upper tangency fits all 85 observations at $0.3788\px$ RMS. More importantly, a front/dorsal-derived ellipse with frozen size, shape, and orientation predicts rear points beyond those used for translation: six anchors predict the other 18 stations at $\SparseRearRMS\px$ RMS, and guarded contiguous rear holdouts yield $\BlockedRearRMS\px$. Guarded spatial holdouts of the strict model yield $\BlockedJointRMS\px$ across all three features. Withholding the entire dorsal boundary is reported separately as a cross-feature sensitivity test. An earlier tapered formulation achieves a lower numerical residual but violates its own ellipse--line tangency condition and is rejected as a physical frustum. A completed comparison of 28 Curiosity images, including 18 uniformly measured whole frames, demonstrates comparable straight segments in geological scenes but does not establish the frequency of the full measured conjunction. The configuration exhibits manufacturing-like geometric regularity supported by predictive tests; composition and provenance are not determined.

\end{abstract}
\vspace{2pt}
\noindent{\small\textbf{Keywords:} Mars; Gale crater; Curiosity; Mastcam-100; cylinder; subpixel metrology; orthogonal regression; image geometry; occlusion; concretion; technosignature assessment.}\par
\vspace{7pt}
\noindent{\color{fine}\rule{\linewidth}{0.5pt}}\par

\section{Introduction}
\label{sec:intro}

A photograph can establish the projected geometry of an object even when its composition, depth, and origin remain unknown. For the Curiosity Mastcam-100 observation considered here, the central question is unusually specific: can the distinct curved boundary at the foreground end, the exposed curvature at the far end, and the upper longitudinal boundary be represented by one common cylindrical geometry? The central evidentiary question is whether that geometry predicts observed boundary points that did not participate in its registration or fit. Determining physical origin is beyond the present scope.

The target occurs in the image identified by \pdsid, acquired on sol 3556 (7 August 2022). The Mastcam-100 instrument is a fixed-focal-length, narrow-field camera with an instantaneous field of view near \(74\,\mu\mathrm{rad}\) per image pixel, making subpixel analysis appropriate for centimeter-scale objects at several meters' range \cite{malin2017,bell2017}. The target is not a full, free-standing object on a uniform background: its lower outline merges with the regolith, and a dark region adjacent to the near end contains a cast-shadow contribution. Consequently, not every visible intensity transition belongs to the solid's perimeter.

We define material boundaries independently of their suitability for any candidate shape. The upper edge is selected at the bright-body/dark-regolith interface; the lower visible contact is classified as possible foreground soil occlusion rather than automatically being treated as a second silhouette generator. Both terminal curves are registered with explicit parameter constraints. The measured feature families are then tested jointly against constant-section and tapered geometries. The image data and executable analyses are retained in the research repository \cite{bogner2026repo}.

This paper reports measurements, predictive validation, uncertainty structure, and the completed geological comparison. To our knowledge, this is the first publicly documented, reproducible quantitative metrology study of this particular feature. Earlier recognition and visualization are acknowledged: Ward credits Rami Bar Ilan with discovery, and MILTOID presented Bayer decoding and a visual 3D cylinder overlay in October 2025 \cite{ward2025,miltoid2025}. Those preliminary checks precede the numerical metrology and predictive tests reported here. The original study appears in the repository history on 24 October 2026 and in the Zenodo deposit published 30 October; independent reconstruction was added on 11 November, followed by the 13 November predictive validation \cite{bogner2026deposit}.

\section{Observational Material and Provenance}
\label{sec:data}

\subsection{Instrument products and pixel conventions}

The primary inputs are NASA PDS3 Reduced Data Record products from the right Mastcam, with product stem \pdsid\ and dataset designation \texttt{MSL-M-MASTCAM-4-RDR-IMG-V1.0}. The investigation operates on zero-based image coordinates \((x,y)\), with \(x\) increasing rightward and \(y\) increasing downward. For the principal radiometric product, an image coordinate refers to the array location in the PDS subframe, not the full detector's absolute position.

\begin{table}[ht]
\centering\small
\caption{Image products and their analytic roles. Dimensions are ordered as image lines \(\times\) samples \(\times\) bands.}
\label{tab:products}
\begin{tabularx}{\linewidth}{@{}l l l X@{}}
\toprule
\textbf{Suffix} & \textbf{Array shape} & \textbf{Encoding} & \textbf{Use in analysis} \\
\midrule
\texttt{DRXX} & \(1184\times1328\times3\) & big-endian 16-bit & Primary radiometric pixel measurements; native CAHVOR geometry \\
\texttt{DRLX} & \(1180\times1323\times3\) & big-endian 16-bit & Independently sampled linearized CAHV geometry \\
\texttt{DRCX} & \(1184\times1328\times3\) & 8-bit & Tonal-display and localization reference \\
\texttt{DRCL} & \(1180\times1323\times3\) & 8-bit & Linearized display and visual comparison \\
\bottomrule
\end{tabularx}
\end{table}

The image metadata give an acquisition time of 2022-08-07 20:58:23.349 UTC, local true solar time 12:10:25, a \(14.2\,\mathrm{ms}\) exposure, and clear-filter Mastcam-right acquisition. The principal data were transmitted with predictive lossless Bayer Huffman compression. Their geometric features therefore are not attributable to JPEG block boundaries. Demosaicing, optical blur, radiometric scaling, and companding nevertheless influence local edge localization and are addressed through multi-band and estimator-sensitivity comparisons.

The per-band radiometric conversion implemented in the repository is
\begin{equation}
I_b(x,y)=o_b+q_b\,\mathrm{DN}_b(x,y),\quad
(q_R,q_G,q_B)=(1.201,1.112,1.037)\times10^{-4},
\label{eq:radiance}
\end{equation}
with \(o_b=0\) in the processed products and radiance in \(\mathrm{W\,m^{-2}\,sr^{-1}\,nm^{-1}}\). The mean of the three band radiances defines the primary scalar field \(I(x,y)\); separate-channel fits provide a check on color interpolation effects.

\subsection{Four distinguishable visual regions}
\label{sec:regions}

The segmentation distinguishes (i) the sunlit longitudinal surface, (ii) its upper background-facing edge, (iii) the foreground curved termination separating illuminated material from the dark recess, and (iv) the lower-right dark patch, whose distal boundary partly follows cast shadow on soil rather than the object's cross-section. The far upper-left termination has its own observable material/background boundary, while its lower part is obscured by regolith.

This classification is causally important. An elliptical fit to a shadow boundary measures illumination geometry and ground topography, not necessarily the perimeter of the object that cast the shadow. Likewise, a line running through the brightest surface is not necessarily the outer silhouette, even if its projected slope is nearly constant.

The analysis reports a regolith-noise estimate \(\sigma_I=7.782\times10^{-3}\) in its radiometric units, a sunlit-object mean near \(9.848\times10^{-2}\), and a dark-core mean near \(1.358\times10^{-2}\), with an edge signal-to-noise estimate of 10.91 for its selected contrast definition. These measurements characterize signal separation without assuming anything about the object's material.

\begin{figure}[ht]
\centering
\includegraphics[width=0.94\linewidth]{figures/target-boundaries.pdf}
\caption{Native DRXX radiometric image, enlarged without generative enhancement. The unmarked crop (left) and the same crop with recorded stations (right) identify the near rim, rear terminus, dorsal material boundary, excluded adjacent-pebble station, and soil/shadow contacts. Display contrast is normalized for visibility; measurements use the original radiometric arrays. Labels for soil and shadow identify interpreted regions rather than newly sampled contours.}
\label{fig:target}
\end{figure}

\section{Camera Geometry, Illumination, and Scale}
\label{sec:camera}

\subsection{CAHV and CAHVOR projection}

Let \(\mat C\) be the camera center, \(\mat A\) its optical axis, and \(\mat H,\mat V\) the horizontal and vertical projection vectors. In the linear CAHV model, a three-dimensional point \(\mat P\) projects to
\begin{equation}
 x=\frac{(\mat P-\mat C)\cdot\mat H}{(\mat P-\mat C)\cdot\mat A},\qquad
 y=\frac{(\mat P-\mat C)\cdot\mat V}{(\mat P-\mat C)\cdot\mat A}.
 \label{eq:cahv}
\end{equation}
The corresponding nonlinear CAHVOR model applies radial corrections about distortion axis \(\mat O\):
\begin{align}
\mat p&=\mat P-\mat C, & \omega&=\mat p\cdot\mat O, \\
\boldsymbol{\lambda}&=\frac{\mat p}{\omega}-\mat O, & \tau&=\boldsymbol{\lambda}\cdot\boldsymbol{\lambda},\\
\mu&=R_0+R_1\tau+R_2\tau^2, &
\mat p'&=\mat p+\mu(\mat p-\omega\mat O).
\end{align}
The corrected vector \(\mat p'\) replaces \(\mat p\) in \cref{eq:cahv}. The script also inverts image coordinates to viewing rays and transforms rover-navigation vectors to the local site reference frame. The independently derived linearized camera azimuth and elevation agree with the recorded pointing values of \(227.7412^\circ\) and \(-22.0047^\circ\), respectively.

At the target location near \((1061,980)\), the singular values of the local angular-to-image Jacobian are \(13{,}498.24\) and \(13{,}418.43\,\mathrm{px\,rad^{-1}}\). Accordingly,
\begin{equation}
 f_{\mathrm{local}}=13{,}458.28\,\mathrm{px\,rad^{-1}},\quad
 \mathrm{IFOV}=74.304\,\mu\mathrm{rad\,px^{-1}},\quad
 \epsilon_{\mathrm{anisotropy}}=0.595\%.
\end{equation}
The small local anisotropy means that a major departure from circularity cannot reasonably be attributed to the lens at this image location.

\subsection{Range-sensitive physical dimensions}

Without an overlapping target stereo solution, angular and pixel-based measurements are better constrained than dimensional measurements in millimeters. Using the metadata's best focus distance \(D_f=3.957\,\mathrm{m}\) as a working range gives a local sampling scale
\begin{equation}
 \ell_{\mathrm{px}}(D)\approx D\,(74.304\times10^{-6})
 = 0.2940\left(\frac{D}{3.957\,\mathrm m}\right)\mathrm{mm\,px^{-1}}.
 \label{eq:scale}
\end{equation}
Within the recorded focus search bracket \(3.6\)--\(4.4\,\mathrm m\), this scale spans approximately \(0.2675\)--\(0.3269\,\mathrm{mm\,px^{-1}}\). A projected strict-model diameter of \(2a\approx39.99\px\) therefore corresponds to \(11.76\,\mathrm{mm}\) at the focus-distance scenario. This is an illustrative physical conversion, not a stereo-derived object range; the true three-dimensional axis length requires an additional orientation estimate.

\subsection{Solar geometry and cast-shadow direction}

PDS metadata give a local solar azimuth of \(186.8169^\circ\) and elevation of \(69.8224^\circ\). Projection through the camera model gives a solar image-plane direction near \(-116.571^\circ\) and a corresponding shadow direction near \(+63.429^\circ\). This direction is consistent with the extensive lower-right dark region. Modeling the shadow as an optical projection onto an underlying soil surface is essential for avoiding an incorrect full-outline fit: the shadow's distal edge need not be concentric with the exposed curved material boundary.

In a simple receiving-plane formulation, let the shadow-casting rim point be \(\mat P\), the unit Sun-direction vector be \(\hat{\mat s}\), and the soil plane satisfy \(\hat{\mat n}_g\cdot(\mat X-\mat P_0)=0\). The physical shadow intersection is
\begin{equation}
 \mat X_{\rm shadow}=\mat P-\lambda\hat{\mat s},\qquad
 \lambda=\frac{\hat{\mat n}_g\cdot(\mat P-\mat P_0)}{\hat{\mat n}_g\cdot\hat{\mat s}},
 \label{eq:shadow}
\end{equation}
when \(\lambda>0\). Camera projection of the resulting point gives the expected image location of the cast shadow. The receiving-plane orientation is a modeled scene property, not an independently measured surface normal. The construction nevertheless explains why the terminal material curve and the lower distal shadow boundary need not be pieces of the same conic.


\section{Subpixel Measurement of Physical Boundaries}
\label{sec:edges}

\subsection{Normal-profile extraction}

For a line segment with unit tangent \(\hat{\mat t}=(t_x,t_y)\), define its directed normal \(\hat{\mat n}=(-t_y,t_x)\). Intensity is sampled on rays
\begin{equation}
 \mat x_i(s)=\mat x_{i,0}+s\hat{\mat n},\qquad
 g_i(s)=I\bigl(\mat x_i(s)\bigr),
\end{equation}
with bilinear interpolation on the radiometric image. After mild Gaussian smoothing, the derivative \(g_i'(s)\) is searched for a contrast extremum of the \emph{physically expected sign}. A local three-sample parabola refines the position; the surrounding unsmoothed profile is also fitted to an error-function transition,
\begin{equation}
 g(s)=I_-+\frac{I_+-I_-}{2}
 \left[1+\operatorname{erf}\left(\frac{s-s_0}{\sqrt 2\,\sigma_{\mathrm{PSF}}}\right)\right],
 \label{eq:erf}
\end{equation}
where \(s_0\) is the transition coordinate, \(\sigma_{\mathrm{PSF}}\) a profile-width parameter, and \(I_\pm\) the fitted asymptotic intensities.

For the upper edge, the selected normal runs from darker background into brighter illuminated object. Its appropriate gradient search polarity is \(+1\). For the lower edge it runs from brighter object into darker soil, requiring \(-1\). The implementation records the sign-adjusted intensity difference across each selected station and checks that the transition follows the expected material boundary.

\subsection{Exterior-boundary validation}
\label{sec:exterior}

The upper boundary is defined by the large, spatially coherent radiometric transition from background regolith into the illuminated body. At each of 26 stations, the extractor records transition contrast, gradient direction, fitted transition position, and fit uncertainty. All 26 stations satisfy the expected polarity diagnostic. Their best-fit line is at $28.895^\circ$, with $0.2422\px$ orthogonal RMS over approximately $36.4\px$ of the visible exterior. The edge-coordinate arrays and diagnostic fields are retained in \sourcepath{step3_features.json}. This feature definition is independent of the subsequent cylinder or tapered-model fits.

\subsection{Near rim: separation from the adjacent pebble}

The near curved termination is sampled at five-degree radial intervals. Of 36 near-rim stations, precisely one at \(\theta=-45^\circ\) lies across a narrow shadow gap on the adjacent upper-right pebble. Its extracted position is approximately \((1083.34,969.16)\), radius \(22.40\px\), compared with neighboring radii of \(19.86\) and \(20.34\px\). The near-rim dataset therefore comprises the 35 stations that lie on the continuous observed termination, with the neighboring pebble excluded by spatial and contrast classification.

\subsection{Far termination: separating cap, soil, and longitudinal transition}

The rear region is sampled from an interior seed at \((1035.5,951.2)\), with radial angles spanning \(-210^\circ\) to \(-35^\circ\). The extractor selects the first outward crossing of a radiometric threshold \(I_{\mathrm{iso}}=0.0860\) before fitting the local transition. The crucial physical partition is between the buried lower-left contact (indices \(0\)--\(7\)), the exposed rear cap (indices \(8\)--\(31\)), and the upper longitudinal transition (indices \(32\)--\(35\)). Only the 24 middle samples are treated as the terminating curve. This prevents a soil-occlusion line from masquerading as a noncircular end.

\section{Curve Fitting and Geometric Models}
\label{sec:models}

\subsection{Orthogonal-distance estimation}

Straightness is estimated using total least squares. For observations \(\{\mat x_i\}\), the eigenvector of the smaller eigenvalue of the 2D sample covariance gives the line normal \(\hat{\mat n}\). Its signed orthogonal residual is
\begin{equation}
 e_i=\hat{\mat n}\cdot(\mat x_i-\bar{\mat x}),\qquad
 \rms_{\perp}=\sqrt{N^{-1}\sum_i e_i^2}.
\end{equation}

For the end curves, an image-plane ellipse centered on \(\mat c\) is parameterized in coordinates aligned with a trial body axis \(\theta\):
\begin{align}
 \hat{\mat u}&=(\cos\theta,\sin\theta), &
 \hat{\mat v}&=(-\sin\theta,\cos\theta),\\
 \mat x(\psi)&=\mat c+b\cos\psi\,\hat{\mat u}
 +a\sin\psi\,\hat{\mat v},
 \label{eq:ellipse}
\end{align}
where \(a\) is the transverse semi-major axis and \(b\) the longitudinally aligned semi-minor axis. A point's residual is its minimum Euclidean distance to the curve, rather than a raw algebraic-conic residual. Newton refinement provides the closest curve point within each nonlinear least-squares evaluation.

\subsection{Uncoupled reference fits}

The 35 clean near-rim observations produce a freely oriented ellipse with \(a=19.7170\px\), \(b=16.4240\px\), \(b/a=0.8330\), and minor-axis angle \(32.704^\circ\), with \(0.16395\px\) RMS. Holding that shape and orientation fixed while translating only its center fits the 24-point rear arc at \(0.18629\px\) RMS. The translated rear center is approximately \((1035.98,964.48)\), relative to the near center \((1068.88,986.00)\), giving a center-to-center bearing of \(33.185^\circ\) and projected separation of \(39.32\px\). This independent two-parameter registration is important because the rear shape is not freely redrawn.

\subsection{Six-parameter strict-cylinder construction}
\label{sec:strict}

The strict model \(M_1\) ties all three observed families to a common, projected constant-section body. Its parameter vector is
\begin{equation}
 \mat p_1=(x_f,y_f,a,b,L,\theta),
\end{equation}
with equal near/far ellipse parameters and collinear centers
\begin{equation}
 \mat c_r=\mat c_f-L\hat{\mat u}(\theta),\qquad
 (a_r,b_r)=(a,b).
 \label{eq:axis}
\end{equation}
The upper image boundary is further constrained to be the outer tangent of that same section:
\begin{equation}
 \hat{\mat v}(\theta)\cdot\bigl(\mat x_d-\mat c_f\bigr)=-a.
 \label{eq:tangent}
\end{equation}
Thus there is \emph{no independently adjustable upper-line offset}. Let \(F\) index the near rim, \(R\) the far arc, and \(D\) the upper boundary. The joint least-squares objective is
\begin{equation}
 \chi^2(\mat p_1)=\sum_{i\in F}\left(\frac{e_{F,i}}{\sigma_{F,i}}\right)^2
 +\sum_{i\in R}\left(\frac{e_{R,i}}{\sigma_{R,i}}\right)^2
 +\sum_{i\in D}\left(\frac{e_{D,i}}{\sigma_{D,i}}\right)^2.
 \label{eq:loss}
\end{equation}
Raw pixel RMS is reported separately as
\begin{equation}
 \rms_{\mathrm{joint}}=\sqrt{\frac{\sum_{i\in F\cup R\cup D}e_i^2}{85}}.
 \label{eq:rms}
\end{equation}
This ensures that differences in the formal uncertainties \(\sigma_i\) do not get silently confused with the reported geometric residuals.

The resulting fit parameters are
\begin{equation}
\begin{aligned}
 (x_f,y_f)&=(1068.3014,985.9459),\qquad
 (x_r,y_r)=(1033.0267,965.3205),\\
 a&=19.9925\px,\quad b=15.7391\px,\quad
 L=40.8622\px,\\
 \theta&=30.3152^\circ,\qquad s=a_r/a_f=1\quad(\text{fixed}).
\end{aligned}
\end{equation}
The fitted projected geometry is highly constrained: a common orientation, equal section size and shape, collinear centers, and an exact upper tangent are enforced simultaneously. In image-space terms, these are precisely the relations expected of a projected cylindrical body. The observational model is a coupled two-dimensional geometric representation, not an independently triangulated 3D surface scan.

\subsection{More flexible comparison models}

Three additional models clarify which geometric constraints are doing work. \(M_2\) retains equal cross-sections and collinearity but estimates a separate upper-crest offset \(r_d\), increasing parameter count from six to seven. \(M_3\) searches projected-axis angle in \(0.1^\circ\) increments, minimizing unweighted residuals while also allowing \(r_d\); because the angle is selected from the data, it has seven effective fitted quantities. \(M_4\) is the earlier eight-parameter tapered image-plane formulation, originally interpreted as a conical frustum, including far-to-near scale ratio \(s\), and uses a tapered upper boundary,
\begin{equation}
 r(u)=r_d\left[1+(s-1)\frac{-u}{L}\right],\quad
 \tan\alpha=\frac{(1-s)r_d}{L},\quad
 \theta_d\approx\theta-\alpha.
 \label{eq:taper}
\end{equation}
Its additional freedom absorbs the observed angular spread numerically, but the formulation fails the physical tangency check in \cref{sec:taper}. \Cref{sec:taper} separates that coupled optimization from an independently translated far-arc comparison.

\section{Results}
\label{sec:results}

\subsection{The joint result across all 85 measured points}

\Cref{tab:joint} summarizes the four joint models. The unweighted raw RMS of strict \(M_1\) is \(0.3788\px\), partitioned into near-rim \(0.2376\px\), far-arc \(0.4830\px\), and upper-edge \(0.4217\px\). Models with an adjustable upper offset yield \(0.3468\px\) (weighted \(M_2\)) or \(0.3389\px\) (unweighted \(M_3\)).

\begin{table}[ht]
\centering\small
\caption{Current joint-fit results from \sourcepath{step9_joint_model.json}; \(N=35+24+26=85\) observations. RMS figures are unweighted orthogonal residuals in native image pixels. Model \(M_3\) counts its scanned angle as a fitted parameter.}
\label{tab:joint}
\begin{tabularx}{\linewidth}{@{}l c c c c c c@{}}
\toprule
\textbf{Fit} & \(k\) & \(s\) & \(\theta\) & \textbf{Near} & \textbf{Far} & \textbf{All} \\
\midrule
\(M_1\), strict tangent & 6 & 1.000 & \(30.32^\circ\) & 0.238 & 0.483 & \textbf{0.379} \\
\(M_2\), offset & 7 & 1.000 & \(31.10^\circ\) & 0.231 & 0.305 & 0.347 \\
\(M_3\), unweighted offset & 7 & 1.000 & \(30.60^\circ\) & 0.242 & 0.378 & 0.339 \\
\(M_4\), invalid frustum & 8 & 0.619 & \(39.33^\circ\) & 0.193 & 0.268 & \textbf{0.255} \\
\bottomrule
\end{tabularx}
\vspace{4pt}
\begin{minipage}{0.96\linewidth}\footnotesize\color{slate}
The upper-edge RMS figures omitted from the compact table are \(0.422\), \(0.485\), \(0.405\), and \(0.310\px\), respectively. The nominal weighted objectives give BIC values of 2680.2, 1947.3, and 1079.6 for \(M_1,M_2,M_4\). These are diagnostics under uncalibrated station errors, not physical-model probabilities. For \(M_1\), \(\chi^2=2653.51\) with 79 nominal residual degrees of freedom gives \(\chi^2_\nu=33.59\). M4 additionally fails its physical silhouette condition.
\end{minipage}
\end{table}



The 6-parameter model's residuals should be understood in relation to its geometric obligations, not only relative to its competitors. It simultaneously puts the rear and near curves on one axis, fixes their common section size, and makes the measured upper exterior edge tangent to that section. Its small, distributed residuals therefore constitute positive evidence for cylindrical consistency. M4's smaller numerical residual does not restore its failed physical tangency constraint.

\subsection{Near- and far-termination agreement}

Independently of the joint fit, the translated unchanged near ellipse matches the far arc at \(0.18629\px\) RMS. The far curve covers only part of the termination; it has no requirement to expose its complete lower perimeter. This registration is therefore evidence that the visible part of the far terminus can carry the same local curvature as the near rim, not an observation of two complete closed end faces.

The test is stronger than a statement that each curved boundary can be approximated by an unrelated ellipse: shape, orientation, and scale are held common. It is weaker than a stereo reconstruction of both terminal planes. The distinction follows directly from the degrees of freedom and the observation geometry.

\subsection{Predictive validation beyond registration points}
\label{sec:prediction}

The new tests preserve the existing coordinates and exclusions. Their protocol was written before these validation runs; the tests are retrospective coordinate-level checks, not a blind repeat of boundary selection. Two front-shape references must be distinguished. The freely oriented ellipse above gives $0.18629\px$ rear RMS. The predictive tests instead freeze the front ellipse constrained to the measured dorsal orientation: $a=20.02897\px$, $b=16.84577\px$, $\theta=28.89475^\circ$, and scale $s=1$. Fitting only its rear translation reproduces $\RegisteredRearRMS\px$ RMS. Neither rear size nor orientation is adjusted.

\begin{table}[H]
\centering\small
\caption{Descriptive fits and heldout predictions. All errors are raw RMS distances in native pixels; the first and fourth rows are in-sample references.}
\label{tab:prediction}
\begin{tabularx}{\linewidth}{@{}l X X r@{}}
\toprule
\textbf{Check} & \textbf{Used in fit} & \textbf{Scored observations} & \textbf{RMS} \\
\midrule
Frozen rear registration & All 24 rear stations; two translations & Same 24 rear stations & \RegisteredRearRMS \\
Sparse registration & Rear indices 0, 4, 8, 12, 16, 20 & Other 18 rear stations & \SparseRearRMS \\
Guarded rear blocks & Remaining rear stations, excluding guards & Each rear station withheld once & \BlockedRearRMS \\
Strict $M_1$, full fit & All 85 stations; six parameters & Same 85 stations & \FullJointRMS \\
Guarded $M_1$ blocks & Remaining stations from all three families & Each of 85 stations withheld once & \BlockedJointRMS \\
Entire dorsal withheld & All front and rear stations only & All 26 dorsal stations & \DorsalWithheldRMS \\
\bottomrule
\end{tabularx}
\end{table}

In five contiguous rear folds, the test block and two neighboring guard stations at each boundary are excluded from translation fitting. The strict-model folds similarly withhold one contiguous block from each feature family, with two-station guards. Every point is scored once across folds. Three prespecified starting angles ($28^\circ$, $31.2^\circ$, $35^\circ$) are tried; the smallest training objective selects the solution, never the withheld score. All starts converge and no parameter reaches a bound. Strict-model heldout RMS values are front $0.8456\px$, rear $0.5946\px$, and dorsal $0.6491\px$; pooling their squared residuals gives $\BlockedJointRMS\px$.

The rear agreement survives both sparse registration and guarded holdouts at approximately $0.2$--$0.24\px$. It therefore extends beyond the rear points used to position the transferred shape. This predicts local curve positions after registration, not the rear's absolute position or cylinder length from the front alone.

\begin{figure}[ht]
\centering
\includegraphics[width=0.95\linewidth]{figures/rear-prediction.pdf}
\caption{Frozen front/dorsal-derived ellipse translated using six rear anchors. Left: native radiometric crop with the transferred curve, fitted anchors, and 18 withheld stations. Right: signed orthogonal residuals for sparse-registration and guarded-block predictions. The curve is clipped to the exposed rear region for legibility; the concealed perimeter is not observed.}
\label{fig:prediction}
\end{figure}

Withholding the entire dorsal feature is a separate sensitivity experiment. Front/rear-only fitting predicts its tangent at $\DorsalWithheldRMS\px$ RMS, maximum $3.391\px$, with axis $33.177^\circ$ rather than the measured dorsal bearing $28.895^\circ$. These modest absolute image deviations are larger than the residuals of the complete fit. They demonstrate the information supplied by the dorsal boundary: removing that measured constraint changes the preferred orientation and tangent. This experiment is not a mandatory passing standard, a measurement of physical surface damage, or a calibrated rejection. It also does not replace the fixed-shape rear tests, which retain the dorsal information by design.

\subsection{Axis alignment and orientation identifiability}

The separate estimates for the near-ellipse minor axis, far-to-near center bearing, and upper line are \(32.704^\circ\), \(33.185^\circ\), and \(28.895^\circ\), respectively. The near ellipse is observed along an incomplete curve, allowing its angle to vary under small changes of the selected points. The reported bootstrap spread for the near-ellipse orientation is approximately \(4.23^\circ\). A transverse shift of \(1\px\) over a \(40.6\px\) center baseline changes its bearing by approximately \(1.41^\circ\). These sensitivities put the several-degree differences in context. Joint optimization chooses \(30.315^\circ\) for the strict cylinder, and achieves subpixel residuals on each point family.

\subsection{Occlusion versus intrinsic taper}

The upper exterior boundary and lower visible diagonal need not represent two opposite silhouette generators. Their measured apparent angular convergence is approximately $35^\circ$, whereas perspective convergence between parallel generators at the measured camera scale is below $0.4^\circ$ for the considered axis inclinations. This excludes the interpretation of \emph{both visible lines} as the unoccluded opposite sides of a parallel-sided body, but does not exclude a partially buried cylinder.

In the partially buried geometry, the upper transition is a material silhouette, while the lower diagonal is the contact between the body and foreground regolith. The latter can cut obliquely across a cylindrical surface as burial depth changes along its length. The two image lines therefore encode different physical boundaries, and their angular separation is not an estimator of cylinder taper without first establishing that both are true generators.

\subsection{Taper: two different statistical questions}
\label{sec:taper}

The rear-arc comparison and the coupled cone fitting answer distinct questions. A \emph{standalone} fit translates an otherwise fixed near-ellipse shape and optionally scales it by \(s\), using only 24 exposed rear points. Its minimum occurs at \(s\approx0.93\) with \(0.18434\px\) RMS, while the exact equal-size case \(s=1\) yields \(0.18629\px\), a difference of \(0.00195\px\). Across \(s=0.85\)--\(1.15\), the variations in raw residual remain small in the context of the image sampling. The conditional profile-likelihood interval recorded by the pipeline is \([0.87,0.99]\); it does not literally include \(s=1\) under that specific weighting model. Raw-fit proximity and formal confidence limits are different statements.

The \emph{coupled} eight-parameter tapered formulation fits all 85 points at once. Its best-fit \(s=0.6187\), with a nominal profile interval \([0.61,0.63]\), is associated with \(\theta\approx39.33^\circ\) and \(\alpha\approx9.46^\circ\). Shrinking the far section allows its proposed silhouette line to rotate toward the observed upper-edge angle, lowering RMS to \(0.2549\px\). A direct physical check rejects the fitted formulation as a frustum silhouette: in axial coordinates its line is $v=-r-mu$, with $m=(1-s)r/L$. Tangency to the end ellipse $u^2/b^2+v^2/a^2=1$ requires $r^2=a^2+m^2b^2$. M4 gives $r=19.294848\px$ whereas the required intercept is $19.433711\px$. Its line therefore cuts its own projected end ellipse. The lower RMS is a reproducible numerical diagnostic, not a valid physical frustum fit. This algebraic failure is distinct from assigning a fixed-front-shape rear-scan deviance to M4, which also changes front shape and orientation.

\begin{table}[ht]
\centering\small
\caption{Two taper evaluations, with different constraints and inferential meaning. Their conditional intervals and objective functions should not be interchanged.}
\label{tab:taper}
\begin{tabularx}{\linewidth}{@{}l X X@{}}
\toprule
 & \textbf{Standalone far-arc registration} & \textbf{Rejected tapered formulation} \\
\midrule
Observed set & 24 exposed rear points & All 85 selected points \\
Free relation & Rear center translates independently & Rear center lies on one fitted axis \\
Best diameter scale \(s\) & 0.930 & 0.619 \\
Nominal 95\% interval & [0.87, 0.99] & [0.61, 0.63] \\
Relevant residual & 0.1843 px rear RMS & 0.2549 px joint RMS \\
Equal-size comparison & 0.1863 px rear RMS at \(s=1\) & Distinct restricted joint model \(M_1\) \\
Interpretation & Partial far arc permits close equal-size registration & Numerical taper violates physical tangency \\
\bottomrule
\end{tabularx}
\end{table}

The measured boundaries admit a stringent constant-section cylindrical configuration. M4 captures the image geometry numerically but fails its own physical tangency condition and is rejected as a physical frustum. This does not test every possible cone or uniquely establish material or origin. A full camera-ray and occlusion-constrained alternative fit remains distinct work.

\section{Sensitivity and Repeatability}
\label{sec:sensitivity}

Measurement stability is assessed not merely by the visual appearance of fitted curves but by changing the estimation rule while keeping the scene fixed. The compiled output compares error-function transition centers, parabolically interpolated gradient maxima, and integer-pixel stations. The two subpixel methods return nearly the same upper-edge angle and consistently smaller line residuals than the discrete alternative (\cref{tab:sens}). An additional trimming experiment removes stations near the segment endpoints; the upper-line orientation shifts moderately while its central portion becomes straighter, as expected if boundary transitions are most ambiguous near neighboring objects and termini.

\begin{table}[ht]
\centering\small
\caption{Sensitivity measurements in the committed \sourcepath{measurements.json}. Values reflect alternative estimators and restricted subsets of the same target, not independent Mars images.}
\label{tab:sens}
\begin{tabularx}{\linewidth}{@{}l c c X@{}}
\toprule
\textbf{Measurement procedure} & \(\theta_{\rm upper}\) & \textbf{Upper-line RMS} & \textbf{Interpretation} \\
\midrule
Error-function subpixel & \(28.895^\circ\) & \(0.2422\px\) & Selected exterior boundary \\
Parabolic gradient peak & \(28.933^\circ\) & \(0.2105\px\) & Close independent localization rule \\
Integer-pixel positions & \(29.105^\circ\) & \(0.4559\px\) & Quantization inflates residual \\
Trim outer 15\% & \(29.154^\circ\) & \(0.1742\px\) & Central line remains straight \\
Trim outer 25\% & \(29.609^\circ\) & \(0.1195\px\) & Fewer stations and shorter baseline \\
\bottomrule
\end{tabularx}
\end{table}

In separate synthetic correlated-noise perturbations, the reported standard deviation of the apparent lower/upper convergence grows from \(1.08^\circ\) at one times empirical noise to \(2.45^\circ\) at three times noise. The method is therefore sensitive to optical sampling and pixel covariance, but the resulting uncertainty can be investigated explicitly rather than hidden behind the assumption of perfect edge localization. These tests characterize the stability of selected image boundaries. Physical tolerance and error-model limits are collected in \cref{sec:limits}.

\section{Texture, Straightness, and Illumination as Independent Constraints}
\label{sec:comparisons}

The selected exterior edge has \(0.2422\px\) orthogonal line RMS over a sampled span of approximately \(36.4\px\). At the focus-distance reference scale that corresponds to an image-plane residual length of about \(71\,\mu\mathrm m\). This is \emph{a projected measurement residual}, not profilometric roughness of a three-dimensional surface at micrometer resolution. A quadratic curvature estimate for the same edge gives a maximum sagitta of \(0.280\px\).

The completed comparison uses 28 retrieved Curiosity images, including 18 whole Mastcam frames measured uniformly and ten captioned geological/context images. A common contour method applies $2\times2$ area averaging, brightness-level contours, and 26 equally spaced samples over a 36.4-native-pixel span. Its target window is selected by proximity to the existing dorsal stations, not minimum RMS. It yields $0.303\px$ RMS, a separate measurement from the full-resolution erf result of $0.2422\px$.

Among 226,426 control windows, 835 ($0.369\%$) are at least as straight; qualifying segments occur in 17 of 18 frames. Straight segments therefore occur in geological scenes. These overlapping photometric windows are correlated and may follow material, fracture, shadow, or shading boundaries; they are not counts of independent rocks. Their fraction is not an origin probability. Smoothing and span change the conditional fraction, so the full sensitivity tables are retained in the supplementary comparison report \cite{comparison2026}.

The conjunction of terminal curvature, transferred geometry, and longitudinal constraint remains the relevant object-level result. The comparison did not estimate the frequency of that complete conjunction. Its automatic rim contour differs from the recorded rim stations by 2.88 pixels on average, preventing a comparable rim ranking; appearance-search scores likewise are not physical-model likelihoods. No full counterpart of the target's 85-station reconstruction was demonstrated by that study. The detailed image inventory, comparison atlas, preprocessing sensitivity, and appearance retrievals remain supplementary material.

Three locally selected natural rock boundaries in the same image provide secondary context. Their reported line RMS values are \(0.971\), \(1.724\), and \(1.953\px\), respectively, compared with \(0.242\px\) for the target (\cref{tab:rocks}). The target edge is therefore approximately four to eight times straighter in this specific local comparison. This describes the three selected local boundaries; the completed comparison above supplies the broader measured context.

\begin{table}[ht]
\centering\small
\caption{Same-image comparative morphology from the current report. Texture is variance of a discrete image Laplacian on normalized radiance patches.}
\label{tab:rocks}
\begin{tabularx}{\linewidth}{@{}l c c c@{}}
\toprule
\textbf{Sample} & \textbf{Edge RMS (px)} & \textbf{Relative edge RMS} & \textbf{Texture variance} \\
\midrule
Target & \textbf{0.242} & 1.0 & \(8.02\times10^{-5}\) \\
Nearby clast 1 & 0.971 & 4.0 & \(1.64\times10^{-4}\) \\
Nearby clast 2 & 1.724 & 7.1 & \(2.45\times10^{-4}\) \\
Nearby clast 3 & 1.953 & 8.1 & \(2.98\times10^{-4}\) \\
\bottomrule
\end{tabularx}
\end{table}

An additional radiometric test examines intensity across the upper body. The current Step 5 output reports curvature coefficient $-0.01285 \pm 0.01357$, with significance $z=0.95$, and relative residuals 0.02682 for the planar fit and 0.02673 for the quadratic fit. Historical prose quoted 7.38 or 11.45 for other or undocumented selections; those values are not reproduced by the current pipeline. The current test does not establish significant curvature or reject planar shading. Albedo, scattering, dust, and illumination geometry further limit shape inference from brightness.

The smoothness, curved terminal boundaries, and tangency-constrained joint fit are complementary constraints. None should be substituted for the others, and the most informative result is their common geometric compatibility.

\section{Geometric Implications and Scope}
\label{sec:discussion}

The evidence is cumulative. The front rim constrains the section; the exposed rear curve tests transferred curvature; and the measured dorsal boundary constrains orientation and exterior tangency. The six-parameter model imposes these relations jointly. Sparse registration and guarded holdouts add evidence beyond a fitted overlay: untrained rear stations continue to follow the frozen shape with approximately $0.2$--$0.24\px$ RMS. A straight sedimentary edge alone is not the identifying result.

Constant-section engineered elements possess the geometric relations tested here. The measured configuration therefore exhibits manufacturing-like geometric regularity, with predictive support rather than visual resemblance alone. The completed geological comparison has not demonstrated a natural analogue reproducing the measured conjunction. Neither naming abrasion, fracture, concretion, or sediment fluidization nor pointing to a straight edge supplies a quantitative explanation of this configuration. Those processes remain documented geological context \cite{bridges2014,loope2011,rubin2017,chan2004,wiens2017}; this study does not require eliminating every conceivable natural process. Manufacture is a morphologically motivated hypothesis, while determining origin is beyond the present scope.

\subsection{Measurement and interpretation limits}
\label{sec:limits}

The validation is retrospective and uses selected coordinates from one acquisition. Neighboring samples share optical and interpolation information, even with coordinate-level guard bands. The nominal reduced chi-square of 33.59 signals an inadequate error/model description; correlation, underestimated localization errors, and model discrepancy are possible contributors, not separately measured explanations. No calibrated origin odds, coverage-calibrated joint confidence interval, or model ranking follows from those nominal weights. Model-specific Birge rescaling is not a common likelihood or a fitted spatial covariance.

The partial terminal arcs constrain visible boundaries rather than both complete hidden end faces. The shared cylindrical model is a constrained image-plane representation, not a unique three-dimensional reconstruction. Focus-based range and receiving-plane assumptions are conditional. Pixel-to-micrometer conversions quantify projected residual distances at an assumed range, not independently measured machining tolerance or surface roughness. The present photometry ($z=0.95$) does not independently establish curved shading.

The comparison study measures photometric features, not a population frequency of the full geometry. Its appearance and straightness rankings cannot be converted into probabilities of natural or manufactured origin. Composition, provenance, and a unique concealed shape remain undetermined. Additional observations could address them, but they are not prerequisites for reporting the measured predictive consistency.


\section{Reproducibility and Data Availability}
\label{sec:reproduce}

The source images, analysis scripts, and numerical outputs are retained in \repo. The original reconstruction base is commit \texttt{90fbf8f89ac2634a9dbd2b27f81fc80d68b9943c} (11 November 2026). The present 13 November revision additionally incorporates the predictive validation, completed comparison, and audit corrections. The old base commit alone does not contain this complete revision.

\noindent\begin{minipage}{\linewidth}
\begin{lstlisting}[language=bash,caption={Reproduction from the repository root; no additional image acquisition is required.}]
python -m pip install -r audit/comparative-study/requirements.txt
python audit/predictive-validation/run_validation.py
python -m unittest discover -s audit/comparative-study/scripts -p test_analysis.py -v
python papers/build_paper_evidence.py
cd papers
xelatex -interaction=nonstopmode -halt-on-error martian-cylinder-technical-paper.tex
xelatex -interaction=nonstopmode -halt-on-error martian-cylinder-technical-paper.tex
\end{lstlisting}
\end{minipage}

The original reconstruction can separately be reproduced with \sourcepath{independent-reconstruction/run_all.py}. Primary outputs include \sourcepath{step3_features.json} (stations), \sourcepath{step8_pebble_and_rear_terminus.json} (front/rear definitions), and \sourcepath{step9_joint_model.json} (joint fits and scale profiles). Predictive folds, fitted parameters, residuals, convergence diagnostics, and input hashes are in \sourcepath{audit/predictive-validation/results.json}. The fixed comparison corpus and its summaries are under \sourcepath{audit/comparative-study/}; no new geological sampling accompanies this revision.

The accompanying \sourcepath{papers/paper-evidence-manifest.json} identifies exact SHA-256 hashes of the manuscript, scripts, numerical inputs, protocols, and comparison records used for the added results. \Cref{app:revision} reproduces the principal hashes. This manifest identifies the working revision without inventing a new committed snapshot or backdating later tests. The figure/data build checks these records and generates the predictive table values directly from JSON.


\section{Conclusion}

Physically directed subpixel measurements yield three mutually constraining features: a 35-point near rim, a 24-point exposed far termination, and a 26-point dorsal boundary. Their strict common-axis, common-section, tangent-boundary model fits all 85 observations at $0.3788\px$ RMS. More decisively, the frozen front/dorsal-derived shape predicts eighteen rear stations after six-anchor registration at $\SparseRearRMS\px$ RMS, and guarded rear holdouts at $\BlockedRearRMS\px$. Guarded spatial prediction across all three feature families yields $\BlockedJointRMS\px$. The tested M4 formulation fails its own physical tangency condition.

The measured configuration exhibits manufacturing-like geometric regularity, supported by frozen-shape transfer and heldout prediction. The completed geological comparison has not demonstrated a natural analogue reproducing the measured conjunction. These findings establish a quantitative morphological anomaly and motivate manufacture as a hypothesis; they do not determine composition or provenance.


\section*{Acknowledgment and provenance note}
\addcontentsline{toc}{section}{Acknowledgment and provenance note}
The calculations and data products synthesized here were developed in the public \emph{martian-cylinder} research repository, using NASA Mars Science Laboratory data. Repository provenance is attributed explicitly; this paper is a synthesis of the recorded investigation, not a claim of independently reacquiring a second rover observation. No NASA institutional endorsement is implied.

\begin{thebibliography}{99}\setlength{\itemsep}{4pt}
\bibitem{bogner2026repo} Bogner, J. (2026). \emph{Martian-cylinder: Independent reconstruction of Curiosity Sol 3556 Mastcam-100 imagery}. Cambrian Minds Media, source code, technical reports, image products, and numerical output. \url{https://github.com/CambrianMinds/martian-cylinder/tree/90fbf8f89ac2634a9dbd2b27f81fc80d68b9943c}.
\bibitem{pds} NASA Planetary Data System and Malin Space Science Systems. \emph{Mars Science Laboratory Mastcam RDR image product \pdsid}; dataset \texttt{MSL-M-MASTCAM-4-RDR-IMG-V1.0}, sol 3556, 7 August 2022. PDS3 labels and image arrays distributed with the repository snapshot.
\bibitem{malin2017} Malin, M. C., et al. (2017). The Mars Science Laboratory (MSL) Mast cameras and Descent Imager: Investigation and instrument descriptions. \emph{Earth and Space Science}, 4, 506--539. \url{https://doi.org/10.1002/2016EA000252}.
\bibitem{bell2017} Bell, J. F., III, et al. (2017). The Mars Science Laboratory Curiosity rover Mastcam instruments: Preflight and in-flight calibration, validation, and data archiving. \emph{Earth and Space Science}. \url{https://doi.org/10.1002/2016EA000219}.
\bibitem{bridges2014} Bridges, N. T., et al. (2014). The rock abrasion record at Gale Crater: Mars Science Laboratory results from Bradbury Landing to Rocknest. \emph{Journal of Geophysical Research: Planets}, 119, 1374--1389. \url{https://doi.org/10.1002/2013JE004579}.
\bibitem{loope2011} Loope, D. B., Kettler, R. M., and Weber, K. A. (2011). Morphologic clues to the origins of iron oxide-cemented spheroids, boxworks, and pipelike concretions, Navajo Sandstone of south-central Utah, U.S.A. \emph{The Journal of Geology}, 119, 505--520. \url{https://doi.org/10.1086/661110}.
\bibitem{rubin2017} Rubin, D. M., et al. (2017). Fluidized-sediment pipes in Gale crater, Mars, and possible Earth analogs. \emph{Geology}, 45(1), 7--10. \url{https://doi.org/10.1130/G38339.1}.
\bibitem{chan2004} Chan, M. A., Beitler, B., Parry, W. T., Orm\"o, J., and Komatsu, G. (2004). A possible terrestrial analogue for haematite concretions on Mars. \emph{Nature}, 429, 731--734. \url{https://doi.org/10.1038/nature02600}.
\bibitem{ward2025} Ward, J. (2025). Cylindrical Object in Colour (Sol 3556), 23 October. \url{https://megalithicmars.com/2025/10/23/rover-curiosity-cylindrical-object-in-colour-sol-3556/}.
\bibitem{miltoid2025} MILTOID (2025). VFX Artist Reacts -- Ancient Artefact on Mars, 23 October. Video; image decoding and visual cylinder overlay. \url{https://www.youtube.com/watch?v=xpr8cWDRTO8}.
\bibitem{bogner2026deposit} Bogner, J. (2026). Morphological and Photometric Analysis of an Unidentified Cylindrical Feature in Curiosity Sol 3556 Mastcam Imagery. Zenodo v1.0.0, published 30 October. \url{https://doi.org/10.5281/zenodo.23789092}.
\bibitem{comparison2026} Bogner, J. (2026). Curiosity geological comparison study, 13 November. Supplementary report, image manifest, and numerical exports, \sourcepath{audit/comparative-study/report.md}, accompanying this paper's evidence manifest.
\bibitem{wiens2017} Wiens, R. C., et al. (2017). Centimeter to decimeter hollow concretions and voids in Gale Crater sediments, Mars. \emph{Icarus}, 289, 144--156. \url{https://doi.org/10.1016/j.icarus.2017.02.003}.
\end{thebibliography}

\appendix
\section{Detailed Measurement Inventory}
\label{app:inventory}

\begin{table}[ht]
\centering\small
\caption{Feature-level source coordinates and elementary fits. All coordinates refer to zero-based DRXX pixels.}
\begin{tabularx}{\linewidth}{@{}l l l X@{}}
\toprule
\textbf{Feature} & \textbf{Count} & \textbf{Observed region} & \textbf{Output / relation} \\
\midrule
Near rim & 35 & \(x:1051.52\)--\(1083.73\); \(y:967.15\)--\(994.17\) & Unconstrained ellipse \(a=19.717\), \(b=16.424\), \(0.164\px\) RMS \\
Far cap & 24 & \(x:1027.91\)--\(1039.23\); \(y:945.40\)--\(949.86\) & Equal-size translated near ellipse, \(0.186\px\) RMS \\
Upper edge & 26 & \(x:1036.84\)--\(1069.03\); \(y:945.40\)--\(962.55\) & ODR line \(28.895^\circ\); \(0.242\px\) RMS \\
Shared fit & 85 & Combined regions & Strict cylinder: \(30.315^\circ\), \(0.379\px\) RMS \\
\bottomrule
\end{tabularx}
\end{table}

The near-rim set retains 35 of 36 sampled radial stations after excluding the station at \(-45^\circ\) that crosses onto the adjacent pebble. The far arc is the bounded exposed portion of a 36-point extraction, excluding both the burial contact and the transition to the upper longitudinal boundary. These criteria are feature-based rather than chosen only by minimizing a fitted residual.

\section{Projection and Identifiability Details}
\label{app:projection}

For an ideal projected circular end plane with normal aligned to the three-dimensional cylinder axis, the end appears elliptical under oblique viewing. In an orthographic approximation its axis ratio satisfies \(b/a\approx|\cos\gamma|\), where \(\gamma\) is the angle between the end-plane normal and the line of sight. The repository's unconstrained near-ellipse ratio \(b/a=0.833\) is therefore compatible with \(\gamma\approx33.6^\circ\), subject to perspective and local projection convention. Its orientation is much less stable than its point-to-ellipse residual because only a partial arc is observed.

The projected position of an unobserved ellipse center is similarly uncertain. Over a short terminal arc, changes in scale and center position can compensate, allowing a one-to-one translation to reproduce the observed curve with only small increases in RMS. This is why rear curvature should be reported as a constraint on the \emph{visible arc}, rather than a direct measurement of the entire concealed end plane. When coupled to a longitudinal surface hypothesis, the same feature gains additional predictive meaning: its center and scale can no longer change independently.

For the strict model, local coordinates of a measured upper-edge point \(\mat x_d=(x_d,y_d)\) are
\begin{align}
 u_d&=(x_d-x_f)\cos\theta+(y_d-y_f)\sin\theta,\\
 v_d&=-(x_d-x_f)\sin\theta+(y_d-y_f)\cos\theta.
\end{align}
The tangent residual \(e_d=v_d+a\) follows immediately from \cref{eq:tangent}. Its sign and the image coordinate orientation are not cosmetic: reversing the orientation of the extracted intensity boundary changes which physical contour is being compared with this equation.

\section{Sensitivity, Resolution, and Error Interpretation}
\label{app:error}

Estimator comparisons and uncertainty-model limitations are discussed in \cref{sec:sensitivity,sec:limits}. The following dimensional scenarios retain the explicit distance dependence of \cref{eq:scale}.

Physical measurements in millimeters depend linearly on assumed distance \(D\). For the strict cylinder, the transverse diameter \(2a\approx39.99\px\) scales to approximately \(10.70\)--\(13.07\,\mathrm{mm}\) across the nominal focus-search interval \(3.6\)--\(4.4\,\mathrm m\), or \(11.76\,\mathrm{mm}\) at \(3.957\,\mathrm m\). The projected center separation \(40.86\px\) is \(12.0\,\mathrm{mm}\) at that same assumed distance; the full three-dimensional object length is not the same quantity and requires orientation reconstruction. Distinguishing projected separation from body length prevents the focus-distance scenario from being mistaken for direct ranging.

\section{File Integrity and Analysis Snapshot}
\label{app:hash}

The provenance manifest records the following SHA-256 digests for primary images. These allow future readers to verify that the same byte-level image arrays are being analyzed.

\noindent{\scriptsize\ttfamily
\begin{tabularx}{\linewidth}{@{}l X@{}}
\toprule
DRXX & 45a3d5e5d32ed823e8a3e2ff0dc0f5ffeb5bd546c81be8d51e06c9f241b5e1fd \\
DRLX & 911e661c99a4453f4b3d4bd13a3d35f049c5df965fcb3c91b014a47c38c7605f \\
DRCX & 8aec700e07d61123151c0811332f76e10f13d1e3fb361de5b10476b4961c473f \\
DRCL & 761c9c1774cfe5152ca7d5cb729de285d08fa78c3df05af8720cbd9599b37592 \\
\bottomrule
\end{tabularx}
}

\medskip
\noindent The original base revision and this paper's evidence manifest serve different roles: the base identifies the reconstruction history; the manifest identifies the exact files used in the current synthesis. Later replacements should be recorded as new revisions.

\section{Evidence Manifest for the 13 November Revision}
\label{app:revision}
\noindent Paths are relative to the repository root. The accompanying JSON manifest lists all synthesis inputs, scripts, and figure digests.
\begin{longtable}{@{}p{0.97\linewidth}@{}}\toprule
\sourcepath{audit/predictive-validation/protocol.md}\\
{\scriptsize\ttfamily 84e2c196bc06eaecdf5dabc5dc0a73cf\allowbreak{}5c21668aba01ff193c01c59943915034}\\[5pt]
\sourcepath{audit/predictive-validation/run_validation.py}\\
{\scriptsize\ttfamily 812ee723153ddc7db95e845771de6e4d\allowbreak{}98d6776ccfdc75c6e9f4681c2bc45699}\\[5pt]
\sourcepath{audit/predictive-validation/results.json}\\
{\scriptsize\ttfamily f3c0edf59268e9a0fde3b6893f3744c0\allowbreak{}f54f0049b07111fac03bd89aee5e3f49}\\[5pt]
\sourcepath{audit/comparative-study/comparison-summary.json}\\
{\scriptsize\ttfamily 82d455289db2d39dba8b66069495b976\allowbreak{}ccbd522aba666dbcf151b9a6040293ff}\\[5pt]
\sourcepath{audit/comparative-study/image-manifest.json}\\
{\scriptsize\ttfamily be6e7f623eea166e647307f409dbe528\allowbreak{}69f1c34576d9de7bbf3548ab3980efda}\\[5pt]
\sourcepath{papers/martian-cylinder-technical-paper.tex}\\
{\scriptsize\ttfamily ae217d343744c0cfc627b75061829fa9\allowbreak{}1b7d48640b82cad791bd310c0458ab7d}\\[5pt]
\bottomrule\end{longtable}

\end{document}
